\chapter*{Giới thiệu}
\addcontentsline{toc}{chapter}{Giới thiệu}

\renewcommand{\arraystretch}{1.6}
\begin{flushleft}
  \begin{tabular}{@{}>{\bfseries\color{deepblue}}l@{\hspace{0.8cm}}p{0.65\textwidth}@{}}
    Sinh viên thực hiện: & Nguyễn Phúc Hậu \\
    Mã số sinh viên: & 23120252 \\
    Email sinh viên: & \textit{23120252@student.hcmus.edu.vn} \\
    Tên môn học: & Kiểm thử phần mềm \\
    Mã lớp: & CQ2023/31 \\
    Giảng viên lý thuyết: & Trần Thị Bích Hạnh \\
    Giảng viên thực hành: & Trương Phước Lộc; Hồ Tuấn Thanh \\
    Năm học: & 2026--2027 \\
  \end{tabular}
\end{flushleft}
\renewcommand{\arraystretch}{1.2}

\section*{Phạm vi bài làm}

Báo cáo trình bày việc áp dụng \textit{Domain Testing} và
\textit{Boundary Value Analysis} trên EShop tại commit
\texttt{85af3ba875c88283615e22cb108f13e2fccaf0e9}. Phạm vi cá nhân gồm bốn
feature: \texttt{A-FR03} quên mật khẩu và đặt lại mật khẩu trên Web,
\texttt{B-FR09} Discount Coupon tại Web checkout, \texttt{C-FR14} quản lý
danh mục trên Admin Web, và \texttt{D-M09} Discount Coupon tại Mobile checkout.

Các ảnh ở đầu mỗi chương chỉ là \textit{smoke evidence}, dùng để xác nhận feature
có thể truy cập trong môi trường kiểm thử. Chúng chưa phải bằng chứng
\textit{Pass/Fail} cho test case chi tiết.

Equivalence Partitioning được dùng để chia miền thành các tập không rỗng, không
chồng lấn; Boundary Value Analysis chọn giá trị tại và quanh biên của các miền
có thứ tự. Cách áp dụng tuân theo ISTQB CTFL v4.0.1, mục 4.2.1--4.2.2
\cite{istqb_ctfl_4_0_1}, đồng thời đáp ứng yêu cầu kỹ thuật của đề HW02
\cite{hw02_assignment_2026}.
